% !TeX root = ../main.tex
% 正确模型下，无成本额外信息的最优决策价值与固定策略表现必须区分。
% 不把估计均方误差排序直接等同于服务决策损失排序。
\chapter{道路信息误差的决策影响与观测配置}
\label{chap:information}

\section{服务韧性分析中的信息约束与研究问题}
\subsection{道路状态与观测信息的区分}
\subsection{固定策略与条件最优决策的比较边界}

\ThesisBlank

\section{道路估计误差向服务损失的传递机制}
\subsection{关键路段时延高估与低估的不同影响}
\subsection{路径切换与容量瓶颈转移的阈值效应}

\ThesisBlank

\section{信息精度与决策价值的关系及成立条件}
\subsection{正确模型下无成本信息的价值边界}
\subsection{任务加权误差与策略损失的排序关系}
\subsection{观测成本与信息到达延迟的净收益影响}

\ThesisBlank

\section{面向服务决策的道路观测配置}
\subsection{观测预算与决策相关目标}
\subsection{一次性配置与序贯更新方法}
\subsection{代理目标的适用条件与可行回退}

\ThesisBlank

\section{误差传递与信息价值的机制实验}
\subsection{相同误差规模下的方向与位置对照}
\subsection{估计精度排序与决策损失排序的差异检验}
\subsection{观测配置与状态更新的收益来源分解}

\ThesisBlank

\section{泸定情景验证与适用边界}
\subsection{不同服务阶段和瓶颈状态下的信息价值}
\subsection{模型失配与误导观测的反例分析}
\subsection{观测成本和更新延迟下的收益边界}

\ThesisBlank

\section{本章小结}
\ThesisBlank
